% ============================================================================
% PHU LUC B: KET QUA TRACE CHAY TUNG BUOC
% ============================================================================

\chapter{KẾT QUẢ CHẠY VÀ DỮ LIỆU ĐẦU RA}
\label{app:trace}

Phụ lục này tổng hợp toàn bộ output chính của stage~1, gồm:
\texttt{tools/trace\_algorithms}, ứng dụng phát hiện port-scan, và dữ liệu
đo hiệu năng/trực quan hóa dùng trong đánh giá. Mục đích là minh họa
\emph{bằng dữ liệu thật} các tính chất lý thuyết đã chứng minh trong
Chương~\ref{ch:compare}:
\begin{itemize}
    \item Count-Min Sketch \emph{ước lượng vượt} ($\hat f(x) \ge f(x)$).
    \item Misra--Gries \emph{ước lượng thiếu} ($\hat f(x) \le f(x)$)
          với sai số bị chặn bởi $m/k$.
    \item Hai giải thuật \emph{kẹp} tần suất thật giữa hai giá trị.
\end{itemize}

\noindent\textbf{Dữ liệu đầu vào.} Luồng
$S = \langle a, b, a, c, a, b, d, a \rangle$, độ dài $m = 8$.
Tần suất chính xác: $f(a)=4$, $f(b)=2$, $f(c)=1$, $f(d)=1$, $f(e)=0$.

\noindent\textbf{Tham số.} Count-Min Sketch với
$w=4, d=2$; Misra--Gries với $k=3$ (tối đa $k-1=2$ bộ đếm).
Kích thước nhỏ được chọn có chủ đích để toàn bộ ma trận và bảng đếm
in được trực tiếp vào báo cáo.

\noindent\textbf{Ghi chú về nguồn dữ liệu.} Các listing trong phụ lục này
được lấy trực tiếp từ output sinh ra khi chạy thực nghiệm trên máy, và
được ghi lại dưới thư mục \texttt{output/} để đảm bảo tái lập.

\section{Menu và chế độ chạy}
\label{sec:trace:menu}

Menu của chương trình \texttt{trace\_algorithms} khi chạy không tham số:

\lstinputlisting[style=pseudo,
    basicstyle=\ttfamily\scriptsize,
    caption={Menu \texttt{trace\_algorithms} (\texttt{output/appendix/trace\_algorithms\_help.txt})},
    label={lst:trace-menu}]{\outputDir/appendix/trace_algorithms_help.txt}

\section{Count-Min Sketch}
\label{sec:trace:cms}

Mỗi bước gồm: chỉ số cột $h_r(x) \bmod w$ trên từng hàng, phép
\texttt{+= 1} tương ứng và toàn bộ ma trận sau khi cập nhật. Phần
\textsc{QUERIES} liệt kê tất cả truy vấn cùng với tập cell được lấy min.

\lstinputlisting[style=pseudo,
    basicstyle=\ttfamily\scriptsize,
    caption={Count-Min Sketch --- trace đầy đủ (\texttt{output/appendix/trace\_algorithms\_cms.txt})},
    label={lst:trace:cms}]{\outputDir/appendix/trace_algorithms_cms.txt}

\paragraph{Nhận xét.}
Ở bước 8, khóa \texttt{d} có $f(d)=1$ nhưng $\hat f(d)=6$. Lý do:
hai cell của \texttt{d} là $(0,0)$ và $(1,0)$ \emph{trùng hoàn toàn}
với hai cell của khóa nóng \texttt{a}. Đây là minh họa cụ thể cho
hiện tượng va chạm hash làm giảm độ chính xác khi $w$ nhỏ. Khi $w$ tăng
theo công thức $w = \lceil e/\varepsilon \rceil$, xác suất xảy ra tình
huống tương tự giảm tuyến tính theo $\varepsilon$.

\section{Misra--Gries}
\label{sec:trace:mg}

Mỗi bước chỉ rõ một trong ba nhánh: (1) khóa đã có trong $T$ nên tăng
bộ đếm, (2) khóa mới và $|T| < k-1$ nên thêm vào, (3) khóa mới và
$|T| = k-1$ nên \emph{giảm toàn bộ} rồi loại các bộ đếm về 0.

\lstinputlisting[style=pseudo,
    basicstyle=\ttfamily\scriptsize,
    caption={Misra--Gries --- trace đầy đủ (\texttt{output/appendix/trace\_algorithms\_misra.txt})},
    label={lst:trace:mg}]{\outputDir/appendix/trace_algorithms_misra.txt}

\paragraph{Nhận xét.}
Hai sự kiện \emph{decrement-all} xảy ra ở bước 4 và bước 7, mỗi sự kiện
loại hai bộ đếm cùng một đơn vị. Với $m=8$ và $k=3$ ta có $m/k = 2$,
trùng đúng với sai số lớn nhất quan sát được trên khóa \texttt{a}
($f=4, \hat f=2$). Điều này khớp với định lý Misra--Gries:
$f(x) - \hat f(x) \le m/k$.

\section{So sánh}
\label{sec:trace:compare}

Bảng dưới đây đặt cạnh nhau kết quả của ba phương pháp: đếm chính xác
(HashMap), Count-Min Sketch và Misra--Gries trên cùng luồng
$S = \langle a, b, a, c, a, b, d, a \rangle$.

\lstinputlisting[style=pseudo,
    basicstyle=\ttfamily\scriptsize,
    caption={So sánh ba phương pháp (\texttt{output/appendix/trace\_algorithms\_compare.txt})},
    label={lst:trace:compare}]{\outputDir/appendix/trace_algorithms_compare.txt}

\paragraph{Kết luận.}
Với mỗi khóa, $\textsc{MG}(x) \le f(x) \le \textsc{CMS}(x)$, tức là
Misra--Gries và Count-Min Sketch tạo thành một \emph{cặp chặn hai phía}
cho tần suất thật. Trong thực tế, việc chạy \emph{đồng thời} hai cấu
trúc trên cùng một luồng dữ liệu với kích thước nhỏ ($O(\log n)$ ô
nhớ) cho phép trả lời hai câu hỏi: (i) khóa này có phải heavy hitter
không? (dùng Misra--Gries), (ii) tần suất ước lượng tối đa của khóa
đó là bao nhiêu? (dùng CMS).

\paragraph{Cách tái tạo.}
Để sinh lại các file output trên, từ thư mục gốc của dự án chạy:

\begin{lstlisting}[style=pseudo, basicstyle=\ttfamily\scriptsize, breaklines=true,
                   breakatwhitespace=false, columns=fullflexible, frame=single]
cmake -S . -B build
cmake --build build --target trace_algorithms

OUT=output/appendix
./build/tools/trace_algorithms          > $OUT/trace_algorithms_help.txt 2>&1
./build/tools/trace_algorithms cms      > $OUT/trace_algorithms_cms.txt
./build/tools/trace_algorithms misra    > $OUT/trace_algorithms_misra.txt
./build/tools/trace_algorithms compare  > $OUT/trace_algorithms_compare.txt
\end{lstlisting}

\noindent Các file được ghi dưới \texttt{output/appendix/} để in trực tiếp vào báo cáo.

\section{Output ứng dụng phát hiện port-scan}
\label{sec:trace:portscan-output}

Phần này liệt kê toàn bộ output của pipeline phát hiện port-scan khi chạy
\texttt{scripts/run\_portscan\_demo.sh}. Input là luồng JSONL tổng hợp,
output gồm trace nội bộ của ba thuật toán, cảnh báo và bảng so sánh.

\subsection{Kịch bản tấn công trên cụm Kubernetes}
\label{sec:trace:k8s-attack}

Trước khi đi vào dữ liệu thô, Hình~\ref{fig:k8s-attack-scenario} minh hoạ
kịch bản tấn công mà luồng \texttt{sample\_stream.jsonl} mô phỏng. Topology
giả lập một cụm Kubernetes nhỏ có 11 pod chia thành bốn subnet --- frontend
(\texttt{10.0.1.X}), backend (\texttt{10.0.2.X}), database/cache
(\texttt{10.0.3.X}) và infra (\texttt{10.0.4.X}) --- cộng thêm một pod attacker
(\texttt{10.0.5.99}) nằm ngoài subnet ứng dụng. Mũi tên xám là lưu lượng nền
hợp lệ giữa các pod theo các \textit{NORMAL\_FLOWS} của
\texttt{generate\_portscan\_stream.py}; mũi tên đỏ nét đứt là 50 gói TCP SYN do
attacker phát ra trong khoảng 5\,ms, mỗi gói tới một cặp
$(\textit{dst\_ip}, \textit{dst\_port})$ ngẫu nhiên trong các subnet ứng dụng.

\begin{figure}[H]
  \centering
  \resizebox{\linewidth}{!}{% Kịch bản tấn công port-scan trên cụm Kubernetes giả lập.
% Include từ Phụ lục B (app-cms-trace.tex).
\begin{tikzpicture}[
    every node/.style={font=\footnotesize},
    pod/.style    ={draw, rectangle, rounded corners=2pt, minimum width=1.9cm,
                    minimum height=0.8cm, align=center},
    fe/.style     ={pod, fill=blue!12},
    be/.style     ={pod, fill=green!15},
    db/.style     ={pod, fill=orange!18},
    infra/.style  ={pod, fill=violet!12},
    attacker/.style={pod, fill=red!25, draw=red!70!black, very thick, font=\footnotesize\bfseries},
    detector/.style={pod, fill=yellow!25, draw=black, dashed, minimum width=2.6cm,
                    minimum height=1.0cm, align=center, font=\footnotesize\itshape},
    normal/.style ={-{Stealth[length=2mm]}, gray!60, thin},
    scan/.style   ={-{Stealth[length=2mm]}, red!75!black, thick, dashed},
    alert/.style  ={-{Stealth[length=2.5mm]}, orange!80!black, very thick},
  ]
  % Subnet boxes
  \node[draw, rounded corners=4pt, dotted, minimum width=8cm, minimum height=2.0cm,
        label={[anchor=west, font=\footnotesize\itshape]above left:Subnet 10.0.1.X (Frontend)}]
       (feNet) at (0, 4.5) {};
  \node[draw, rounded corners=4pt, dotted, minimum width=8cm, minimum height=2.0cm,
        label={[anchor=west, font=\footnotesize\itshape]above left:Subnet 10.0.2.X (Backend)}]
       (beNet) at (0, 2) {};
  \node[draw, rounded corners=4pt, dotted, minimum width=8cm, minimum height=2.0cm,
        label={[anchor=west, font=\footnotesize\itshape]above left:Subnet 10.0.3.X (Database / Cache)}]
       (dbNet) at (0, -0.5) {};
  \node[draw, rounded corners=4pt, dotted, minimum width=8cm, minimum height=2.0cm,
        label={[anchor=west, font=\footnotesize\itshape]above left:Subnet 10.0.4.X (Infra)}]
       (inNet) at (0, -3) {};

  % Frontend pods
  \node[fe] (fe1) at (-2.7, 4.5) {frontend-1\\\texttt{10.0.1.10}};
  \node[fe] (fe2) at (0,    4.5) {frontend-2\\\texttt{10.0.1.11}};
  \node[fe] (fe3) at (2.7,  4.5) {frontend-3\\\texttt{10.0.1.12}};

  % Backend pods
  \node[be] (be1) at (-2.7, 2) {backend-1\\\texttt{10.0.2.10}};
  \node[be] (be2) at (0,    2) {backend-2\\\texttt{10.0.2.11}};
  \node[be] (be3) at (2.7,  2) {backend-3\\\texttt{10.0.2.12}};

  % DB pods
  \node[db] (pg)    at (-1.5, -0.5) {postgres\\\texttt{10.0.3.10}};
  \node[db] (redis) at ( 1.5, -0.5) {redis\\\texttt{10.0.3.11}};

  % Infra pods
  \node[infra] (prom) at (-1.5, -3) {prometheus\\\texttt{10.0.4.10}};
  \node[infra] (graf) at ( 1.5, -3) {grafana\\\texttt{10.0.4.11}};

  % Attacker (separate, off to the side)
  \node[attacker] (att) at (-7, 1.5) {attacker\\\texttt{10.0.5.99}};

  % Detector (right side)
  \node[detector] (det) at (7.5, 1.5) {portscan-pipeline\\(CMS + MG + HashMap)};
  \node[font=\footnotesize\bfseries, fill=red!15, draw=red!70!black, rounded corners=2pt,
        inner sep=3pt, anchor=north, align=center]
        (alert) at ($(det.south) + (0, -0.4)$)
        {\texttt{alerts.json}\\\textit{port\_scan: 10.0.5.99}};

  % Normal traffic (a few representative arrows)
  \draw[normal] (fe1) -- (be1);
  \draw[normal] (fe2) -- (be2);
  \draw[normal] (fe3) -- (be3);
  \draw[normal] (be1) -- (pg);
  \draw[normal] (be2) -- (pg);
  \draw[normal] (be1) -- (redis);
  \draw[normal] (prom.north) -- (be1.south);
  \draw[normal] (graf) -- (prom);

  % Port-scan attack: many edges from attacker to random targets
  \draw[scan] (att) -- (fe1);
  \draw[scan] (att) -- (fe2);
  \draw[scan] (att) -- (fe3);
  \draw[scan] (att) -- (be1);
  \draw[scan] (att) -- (be2);
  \draw[scan] (att) -- (be3);
  \draw[scan] (att) -- (pg);
  \draw[scan] (att) -- (redis);
  \draw[scan] (att) -- (prom);
  \draw[scan] (att) -- (graf);

  % Detection pipeline: tap on subnet -> detector
  \draw[alert, dotted] (3.5, 1.5) -- node[above, font=\footnotesize\itshape] {tcpdump SYN} (det.west);
  \draw[alert] (det.south) -- (alert.north);

  % Legend
  \node[font=\scriptsize\itshape, anchor=north west, align=left,
        fill=white, draw=gray!50, rounded corners=2pt, inner sep=4pt]
        at (-7.5, -2.5)
        {\textbf{Chú thích:}\\
         \textcolor{gray!60}{$\rightarrow$} lưu lượng nền hợp lệ\\
         \textcolor{red!75!black}{$\dashrightarrow$} gói SYN từ attacker\\
         \textcolor{orange!80!black}{$\rightarrow$} luồng phát hiện};
\end{tikzpicture}
}
  \caption{Kịch bản tấn công port-scan trên cụm Kubernetes giả lập.}
  \label{fig:k8s-attack-scenario}
\end{figure}

\noindent Pipeline phát hiện đứng ở phía bên phải, tap vào lưu lượng SYN qua
\texttt{tcpdump}, đẩy qua parser thành \texttt{NetworkEvent}, rồi cho ba bộ
ước lượng tần suất xử lý song song. Khi tần suất ước lượng của
\texttt{src\_ip = 10.0.5.99} vượt ngưỡng $\tau = 40$ trong cửa sổ tumbling
60\,s, một bản ghi cảnh báo được phát ra ở \texttt{alerts.json}. Sự kiện đầu
tiên kích hoạt cảnh báo đúng tại sự kiện thứ 141 của luồng (50 gói background
trước attack + 50 gói background xen kẽ + 41 gói scan đầu tiên), trùng đúng
với phân tích thực nghiệm ở Mục~\ref{sec:portscan-impl} của Chương~\ref{ch:eval}.

Đặc trưng \emph{một-nhiều} của port-scan thấy rất rõ trong hình: từ một IP
nguồn duy nhất (attacker) toả ra nhiều mũi tên đỏ tới hầu hết các pod hợp lệ
trong cụm, trong khi mỗi pod nền chỉ có một vài mũi tên xám tới các đối tác
được phép theo kiến trúc dịch vụ. Đây cũng là lý do detector chọn khoá là
\texttt{src\_ip} chứ không phải cặp \textit{(IP, cổng)}: dù attacker đổi đích
hay đổi cổng, mọi gói vẫn rơi vào cùng một bộ đếm.

\subsection{Menu và log chạy end-to-end}
\label{sec:trace:portscan-menu}

\lstinputlisting[style=pseudo,
    basicstyle=\ttfamily\scriptsize,
    caption={Menu \texttt{portscan-pipeline} (\texttt{output/appendix/portscan\_pipeline\_help.txt})},
    label={lst:portscan-help}]{\outputDir/appendix/portscan_pipeline_help.txt}

\noindent Menu liệt kê tất cả tham số CLI: chọn thuật toán
(\texttt{--algo cms|mg|hashmap|all}), trỏ tới file đầu vào
(\texttt{--input PATH}), đường ra (\texttt{--trace-dir},
\texttt{--alerts}, \texttt{--comparison}), và các siêu tham số của từng cấu
trúc dữ liệu. Tham số mặc định khớp với cấu hình trong
Bảng~\ref{tab:portscan-params}.

\lstinputlisting[style=pseudo,
    basicstyle=\ttfamily\scriptsize,
    caption={Log chạy demo end-to-end (\texttt{output/appendix/portscan\_demo\_stdout.txt})},
    label={lst:portscan-demo-stdout}]{\outputDir/appendix/portscan_demo_stdout.txt}

\noindent Log lần lượt báo: số dòng đầu vào đọc được, kết quả mỗi thuật toán
(số sự kiện đã xử lý, số cảnh báo, bộ nhớ đỉnh, thông lượng), nội dung từng
cảnh báo (IP nguồn, tần suất ước lượng, chỉ số sự kiện), và cuối cùng đường
dẫn các file output. Cả ba thuật toán cùng phát một cảnh báo cho
\texttt{10.0.5.99} tại sự kiện \texttt{\#141} với
$\hat f = 41$ --- đúng ngưỡng $+1$ --- xác nhận tính nhất quán giữa CMS, MG
và HashMap trên cùng đầu vào.

\subsection{Luồng sự kiện đầu vào}
\label{sec:trace:portscan-stream}

\lstinputlisting[style=pseudo,
    basicstyle=\ttfamily\scriptsize,
    caption={Luồng sự kiện tổng hợp (\texttt{output/portscan/sample\_stream.jsonl})},
    label={lst:portscan-stream}]{\outputDir/portscan/sample_stream.jsonl}

\noindent File JSON Lines này là \emph{ground truth} cho toàn bộ phụ lục.
Mỗi dòng là một sự kiện kết nối với bốn trường \texttt{src\_ip},
\texttt{dst\_ip}, \texttt{dst\_port}, \texttt{timestamp\_ns}. 100 dòng đầu là
lưu lượng nền (mỗi dòng cách nhau 1\,ms), 50 dòng tiếp theo là gói SYN từ
attacker (mỗi dòng cách nhau 0{,}1\,ms --- mật độ gấp 10 lần lưu lượng nền,
chính là điểm signature mà CMS phát hiện), và 50 dòng cuối là lưu lượng nền
sau attack. Hạt giống cố định \texttt{seed = 42} đảm bảo file này tái tạo
được bit-identical bằng \texttt{python3 scripts/generate\_portscan\_stream.py
--seed 42 --count 200 --attack-start 100 --scan-probes 50}.

\subsection{Cảnh báo và bảng so sánh}
\label{sec:trace:portscan-alerts}

\lstinputlisting[style=pseudo,
    basicstyle=\ttfamily\scriptsize,
    caption={Danh sách cảnh báo (\texttt{output/portscan/alerts.json})},
    label={lst:portscan-alerts}]{\outputDir/portscan/alerts.json}

\noindent Mảng JSON gồm ba bản ghi --- một cho mỗi thuật toán --- tất cả đều
chỉ về cùng một IP \texttt{10.0.5.99}, cùng tần suất 41, cùng nhãn thời gian
và cùng \texttt{event\_index = 141}. Trường \texttt{window\_id} là 0 vì đợt
quét hoàn tất trong vòng 5\,ms nên nằm trọn trong cửa sổ tumbling đầu tiên
(60\,s). Sự đồng thuận tuyệt đối giữa CMS, MG và HashMap chứng minh rằng
trên luồng có $F_0$ nhỏ và phân biệt rõ heavy hitter, cả ba cấu trúc đều
phát hiện được mà không có false positive nào. Tình huống \emph{khác biệt}
chỉ bộc lộ khi $F_0$ tăng lớn --- khi đó HashMap mở rộng tuyến tính theo
$F_0$ trong khi CMS giữ nguyên 80\,KB (xem Bảng~\ref{tab:throughput} ở
Chương~\ref{ch:eval}).

\lstinputlisting[style=pseudo,
    basicstyle=\ttfamily\scriptsize,
    caption={Bảng so sánh 3 thuật toán (\texttt{output/portscan/comparison.csv})},
    label={lst:portscan-comparison}]{\outputDir/portscan/comparison.csv}

\noindent CSV này là nguồn gốc cho Bảng~\ref{tab:portscan-comparison} ở
Chương~\ref{ch:eval}. Đáng chú ý: với $F_0 = 11$ IP phân biệt, HashMap chỉ
chiếm $\sim$0{,}40\,KB và MG $\sim$0{,}42\,KB, còn CMS giữ cố định
$2048 \times 5 \times 8 = 81\,920$ byte $\approx 80$\,KB --- chi phí cố định
mà CMS phải trả để có được tính chất \emph{mergeable} và \emph{độc lập với
$F_0$}, hai tính chất quyết định ở quy mô vận hành thật.

\subsection{Trace nội bộ của ba bộ ước lượng}
\label{sec:trace:portscan-internal}

Ba listing dưới đây là \emph{cốt lõi} của phụ lục --- cho thấy chính xác
trạng thái nội tại của từng cấu trúc dữ liệu tại các mốc sự kiện
50, 100, 110, 120, 130, 140 (trước attack, đầu attack, giữa attack, ngay
trước alert).

\lstinputlisting[style=pseudo,
    basicstyle=\ttfamily\scriptsize,
    caption={Trace CMS trên luồng port-scan (\texttt{output/portscan/trace\_cms.txt})},
    label={lst:portscan-trace-cms-full}]{\outputDir/portscan/trace_cms.txt}

\noindent CMS trace in cấu hình $w \times d$, dung lượng bộ nhớ và watchlist
gồm 8 IP đại diện (attacker + 7 nguồn nền). Chú ý cột \texttt{est} của
\texttt{10.0.5.99}: $0 \to 0 \to 10 \to 20 \to 30 \to 40$ tăng tuyến tính
sau khi attacker bắt đầu phát SYN ở sự kiện 100. Các nguồn nền giữ ổn định
quanh mốc 10--19 --- không có hiện tượng va chạm CMS làm vỡ ước lượng vì
xác suất va chạm $1/w = 1/2048$ rất nhỏ so với $F_0 = 11$.

\lstinputlisting[style=pseudo,
    basicstyle=\ttfamily\scriptsize,
    caption={Trace Misra--Gries trên luồng port-scan (\texttt{output/portscan/trace\_mg.txt})},
    label={lst:portscan-trace-mg-full}]{\outputDir/portscan/trace_mg.txt}

\noindent MG trace bổ sung mục \texttt{top slots} liệt kê tối đa 8 slot có
counter cao nhất (đã giải mã ngược 4 byte little-endian thành dạng IPv4).
Quan sát: với $k = 64$ slot mà chỉ có 11 IP phân biệt, không xảy ra
\emph{decrement-all}, do đó counter của attacker bằng đúng tần suất thật
$f = 41$. Nếu $k$ giảm xuống dưới $F_0$, decrement-all sẽ bắt đầu xảy ra,
counter MG sẽ là \emph{ước lượng thiếu} so với HashMap, và alert có thể trễ
hơn hoặc bỏ sót.

\lstinputlisting[style=pseudo,
    basicstyle=\ttfamily\scriptsize,
    caption={Trace HashMap trên luồng port-scan (\texttt{output/portscan/trace\_hashmap.txt})},
    label={lst:portscan-trace-hashmap-full}]{\outputDir/portscan/trace_hashmap.txt}

\noindent HashMap là baseline chính xác --- counter trong watchlist trùng
đúng với tần suất thật của từng IP trong cửa sổ. So khớp ba trace ta thấy:
trên luồng nhỏ thuần khiết này, cả ba thuật toán cho cùng một con số tại
mọi mốc, sự khác biệt chỉ thể hiện ở \emph{chi phí} (bộ nhớ và thông lượng,
xem Mục~\ref{sec:trace:portscan-alerts}).

\section{Output đo hiệu năng và dữ liệu trực quan hóa}
\label{sec:trace:bench-output}

Kết quả đo hiệu năng và dữ liệu dùng để vẽ biểu đồ được sinh bởi
\texttt{scripts/plot\_bench.py}. Các file \texttt{.csv} được nhúng thành
bảng trong Chương~\ref{ch:eval}; các file \texttt{.dat} là đầu vào trực tiếp
cho gói \texttt{pgfplots} nếu cần vẽ lại biểu đồ.

\lstinputlisting[style=pseudo,
    basicstyle=\ttfamily\scriptsize,
    caption={Thông lượng và bộ nhớ (\texttt{output/bench/throughput.csv})},
    label={lst:bench-throughput-csv}]{\outputDir/bench/throughput.csv}

\noindent Mỗi dòng là kết quả của một benchmark Google Benchmark, gồm tên
benchmark (CMS với $w$ khác nhau, HashMap, hoặc estimate), số iter, thời gian
trung bình mỗi op (ns), thông lượng (ops/s) và bộ nhớ thực
(\texttt{memory\_bytes}). Đây là nguồn gốc của Bảng~\ref{tab:throughput} ở
Chương~\ref{ch:eval}.

\lstinputlisting[style=pseudo,
    basicstyle=\ttfamily\scriptsize,
    caption={Sai số CMS theo chiều rộng (\texttt{output/bench/error\_rate.csv})},
    label={lst:bench-error-csv}]{\outputDir/bench/error_rate.csv}

\noindent CSV này gồm năm cột: chiều rộng $w$, $\varepsilon \approx e/w$, cận
lý thuyết $\varepsilon m$, sai số tuyệt đối đo được trên $10^5$ khoá, và tỉ lệ
\textit{đo / cận}. Tỉ lệ ổn định $28$--$35\%$ qua mọi $w$ là minh chứng thực
nghiệm cho việc bất đẳng thức Markov dùng trong chứng minh CMS là lỏng.

\lstinputlisting[style=pseudo,
    basicstyle=\ttfamily\scriptsize,
    caption={Dữ liệu vẽ throughput (\texttt{output/bench/throughput.dat})},
    label={lst:bench-throughput-dat}]{\outputDir/bench/throughput.dat}

\lstinputlisting[style=pseudo,
    basicstyle=\ttfamily\scriptsize,
    caption={Dữ liệu vẽ sai số (\texttt{output/bench/error\_rate.dat})},
    label={lst:bench-error-dat}]{\outputDir/bench/error_rate.dat}

\noindent Hai file \texttt{.dat} là format đơn giản dạng \emph{x y} (cách
nhau bằng tab), tương thích trực tiếp với \texttt{\textbackslash addplot
table\{...\}} của \texttt{pgfplots} --- người đọc có thể tái tạo
Hình~\ref{fig:bench-throughput} và Hình~\ref{fig:bench-error} ở
Chương~\ref{ch:impl} chỉ bằng cách trỏ \texttt{\textbackslash input} tới
hai file này.

\lstinputlisting[style=pseudo,
    basicstyle=\ttfamily\scriptsize,
    caption={Kết quả đo hiệu năng thô (\texttt{output/bench/bench\_raw.json})},
    label={lst:bench-raw-json}]{\outputDir/bench/bench_raw.json}

\noindent JSON output gốc của Google Benchmark, gồm metadata máy đo
(CPU, kernel, build type) và mảng \texttt{benchmarks} với toàn bộ trường gốc
(real\_time, cpu\_time, items\_per\_second, custom counters). File này được
giữ nguyên để đảm bảo \emph{tính tái lập} --- bất kỳ ai cũng có thể chạy lại
\texttt{bench\_frequency} trên máy mình rồi diff để xem khác biệt phần cứng
ảnh hưởng tới throughput thế nào.
